% SOURCE_AUTHORITY: sga5_fr_workpass.tex, lines 12000--12044 (LF-normalized unit).
% SOURCE_SHA256: 791F4EFFC5E02832D5D77ED03518C8156D6F07E4C8238B03545DB93D883FBB28
\subsection*{8. Tor de complexos}

Sejam $A$ um anel não necessariamente comutativo e ${}_A\cC$ (resp. $\cC_A$) a categoria abeliana dos $A$-módulos à esquerda (resp. à direita). Denotaremos por $K^-({}_A\cC)$ a categoria triangulada dos complexos de ${}_A\cC$ limitados superiormente, e por ${}_A\cD$ a subcategoria triangulada plena de $K^-({}_A\cC)$ dos complexos cujas componentes são $A$-módulos planos à esquerda.

Seja $X^\bullet$ um objeto de $K^-(\cC_A)$, e consideremos o funtor
\[
\Phi_{X^\bullet}=X^\bullet\otimes_A\,\cdot\;:\;K^-({}_A\cC)\longrightarrow K^-(\Ab).
\]
Nestas condições, são válidas as seguintes afirmações:

1) Todo objeto acíclico de ${}_A\cD$ é transformado por $\Phi_{X^\bullet}$ em um objeto acíclico de $K^-(\Ab)$. De fato, sendo $N^\bullet$ um objeto de ${}_A\cD$ tal que $Z^i(N^\bullet)$, $B^i(N^\bullet)$, $H^i(N^\bullet)$ são planos, temos o isomorfismo clássico:
\[
\bigoplus_{p+q=n}H_p(X^\bullet)\otimes H_q(N^\bullet)
\longrightarrow H_n(X^\bullet\otimes_A N^\bullet).
\]

2) Todo objeto $Y^\bullet$ de $K^-({}_A\cC)$ é o alvo de um quase-isomorfismo cuja fonte é um objeto de ${}_A\cD$. Basta, para isso, tomar uma resolução de Cartan--Eilenberg. Podemos, pois, aplicar o teorema 2 do §2, n.º 2 de [2] para concluir que $\Phi_{X^\bullet}$ admite um funtor derivado à esquerda:
\[
\mathbf L\Phi_{X^\bullet}:D^-({}_A\cC)\longrightarrow D^-(\Ab).
\]
O objeto $\mathbf L\Phi_{X^\bullet}(Y^\bullet)$ pode ser obtido, a menos de isomorfismo, da maneira seguinte. Sejam $P^\bullet$ um objeto de ${}_A\cD$ e $P^\bullet\to Y^\bullet$ um quase-isomorfismo. Então $\mathbf L\Phi_{X^\bullet}(Y^\bullet)$ é isomorfo a
\[
Q(X^\bullet\otimes P^\bullet),
\]
onde se denotou, como de costume, por $Q$ o funtor canônico
\[
K^-(\Ab)\longrightarrow D^-(\Ab).
\]
Resulta, portanto, que, para cada objeto $Y^\bullet$ de $D^-({}_A\cC)$, o funtor
\[
K^-(\cC_A)\longrightarrow D^-(\Ab),\qquad
X^\bullet\longmapsto \underline{\underline L}^{-}\Phi_{X^\bullet}(Y^\bullet)
\]
é exato e se anula sobre os complexos acíclicos, do que resulta um bifuntor exato
\[
D^-(\cC_A)\times D^-({}_A\cC)\longrightarrow D^-(\Ab),
\]
\[
(X^\bullet,Y^\bullet)\longmapsto X^\bullet\overset{\mathbb L}{\otimes}_A Y^\bullet.
\]
Passando à homologia, definem-se os $\operatorname{Tor}_i^A(X^\bullet,Y^\bullet)$ (cf. também [1], 6.3). No caso particular de um anel comutativo $A$, obtém-se, de fato, um bifuntor exato
\[
D^-(\cC)\times D^-(\cC)\longrightarrow D^-(\cC),
\]
o produto tensorial no sentido das categorias derivadas ($\cC=\cC_A={}_A\cC$).
